% !TEX root = ../notes_missingsemester.tex
% Block 2: Version Everything — Git & GitHub
%%%%%%%%%%%%%%%%%%%%%%%%%%%%%%%%%%%%%%%%%%%%%%%%%%%%%%%%%%%%%%%%%%%%%%

\index{Git}\index{GitHub}\index{version control}
\chapter{Version Everything: Git \& GitHub}
\newthought{Git a Life.}

% ==============================================================================
% SESSION A: THEORY
% ==============================================================================

% -----------------------------------------------------------------------------
\section{The Why}


\marginnote{\newthought{PixelWise} needs to live on two machines.
Right now there is no reliable way to move code from \texttt{dev} to \texttt{prod}
and keep both in sync. Git gives us exactly that: a shared history
that either machine can pull from, so deployment becomes a deliberate, testable, repeatable step
instead of manual file copying.}

Block~1 gave us two machines and a way to reach one from the other.
We can write code on \texttt{dev} and, eventually, deploy it to \texttt{prod}.
But nothing yet records what changed between deployments, who changed it, or how to undo a mistake.
The moment two people touch the same file, the same function,
or one experiment needs to coexist with another,
the workflow breaks down entirely.

\newthought{Version control}\index{version control} solves an entire family of problems at once.
It records every change with a message, a timestamp, and an author,
so the full history of a project is always available.
It lets you return to any previous state in seconds,
turning risky experiments into cheap, reversible ones.

\newthought{Parallel development} becomes natural.
Multiple people work on the same codebase without overwriting each other's changes.
A new feature, a bug fix, and a refactor can all proceed simultaneously on isolated branches
and be integrated later with explicit, reviewable merge steps.

\newthought{Downtime shrinks} as a direct consequence.
Developers merge through a guided review step that catches conflicts and errors
before they reach \texttt{main}, deploying only code that has already been integrated and verified.
When something still slips through, you can roll back to the last known-good state in a single command
instead of debugging under pressure while users wait.

\newthought{Combined with a hosting platform} like GitHub\marginnote{GitHub: \url{https://github.com}},
version control becomes the backbone of code review,
continuous integration, issue tracking, and project management, built around 
the software development process.




% -----------------------------------------------------------------------------
\subsection{Hands On Experience}

\newthought{Consider this scenario:} you and a classmate are both improving the PixelWise classifier.
You rewrite the data loader; your classmate changes the model architecture.
You each work on your own laptop for a few hours, then try to combine the results.
Without a system for tracking changes, combining means copying files back and forth,
comparing line by line, and hoping nothing was lost. Worse still, if both of
you would have worked on the same machine, overwriting each others files as you go.

\marginnote{\newthought{Every team} that ships software uses version control.
It is not an optional workflow preference; it is the infrastructure
that makes collaborative, iterative development possible at any scale.
As projects grows, interface contracts\index{interface contracts}, clear agreements 
on inputs, outputs,
and responsibilities between components, will become increasingly important.
Version control lays the foundation by making every change visible and reviewable,
so contracts can be enforced rather than assumed.}

\newthought{Git is that system.} It tracks exactly what changed in every file, 
when, and by whom.
It lets you branch off to try an idea, or build a feature without disturbing stable code,
merge branches back together with full visibility into conflicts,
and review each other's contributions before they reach the shared codebase.
The design principle is simple: never lose work, and never wonder what happened.

\newthought{This second block} introduces Git and GitHub as the version control layer for PixelWise.
By the end you will have
\begin{itemize}
    \item understood the Git data model and why it makes history tamper-proof,
    \item practised the staging workflow of editing, staging, and committing changes,
    \item created branches, merged them, and resolved conflicts,
    \item pushed code to a remote repository on GitHub and collaborated through pull requests,
    \item and configured SSH-based authentication so Git operations are seamless from both VMs.
\end{itemize}


\newpage
% -----------------------------------------------------------------------------
\section{The Git Data Model}\index{Git!data model}

\newthought{Git is not a backup tool.} It is a content-addressable 
filesystem with a version history layered on top.
Understanding the data model makes every command intuitive.

\marginnote{Git was created by Linus Torvalds in 2005 to manage the Linux kernel source code.
It is now the de facto standard for version control in software development.\\
Git: \url{https://git-scm.com/}\\
Pro Git book (free): \url{https://git-scm.com/book/en/v2}}

\begin{description}
    \item[Blob]\index{Git!blob} A file's contents, stored by its SHA-1 hash. The name is irrelevant; only the content matters.
    \item[Tree]\index{Git!tree} A directory listing: maps filenames to blobs (files) or other trees (subdirectories).
    \item[Commit]\index{Git!commit} A snapshot of the entire project at a point in time. Contains: a pointer to the root tree,
    the author, a timestamp, a message, and a pointer to the parent commit(s).
    \item[Branch]\index{Git!branch} A lightweight, movable pointer to a commit. \texttt{main} is a branch.
    Creating a branch is instantaneous; it just creates a new pointer.
    \item[HEAD]\index{Git!HEAD} A pointer to the branch you are currently on. When you commit, \texttt{HEAD} moves forward.
\end{description}

\newthought{Every commit} is identified by a SHA-1 hash, a 40-character hexadecimal string like
\texttt{a3f2b7c...}. This hash is computed from the commit's contents. If anything changes, even a single
character in a single file, the hash changes. This makes Git history tamper-proof.

\newthought{Data hygiene}\index{data hygiene} follows directly from the data model.
Because Git stores every version of every blob permanently, committing a large dataset
means carrying it in the repository history forever, even if you delete the file later.
A 500\,MB training set committed once, modified twice, and then removed still occupies 1.5\,GB of history.
Repositories should contain code, configuration, and small metadata files.
Large or frequently changing binary assets, datasets, and model weights do not belong in Git.

\marginnote{\newthought{Dedicated tools} fill the gap.
Git LFS\index{Git!LFS} replaces large files with lightweight pointers.
Hugging Face Hub\index{Hugging Face} hosts datasets and model weights with built-in Git LFS.
DVC\index{DVC} tracks data pipelines and links datasets to experiments.
Weights \& Biases\index{Weights \& Biases} versions datasets alongside experiment logs.\\
Git LFS: \url{https://git-lfs.com}\\
Hugging Face Hub: \url{https://huggingface.co}\\
DVC: \url{https://dvc.org}\\
W\&B: \url{https://wandb.ai}}

\vspace{1.5em}
The \texttt{.gitignore} file, covered in its own subsection below, is the practical mechanism
for keeping these artefacts out of the repository.





% -----------------------------------------------------------------------------
\newpage
\section{The Staging Workflow}\index{Git!staging}

\marginnote{\newthought{The napkin cycle} as easy as it gets:\\
\texttt{git pull}\\
\texttt{git add <files>}\\
\texttt{git commit -m "..."}\\
\texttt{git push}\\
Pull the latest changes, stage your work, commit it, push it. Every other command is a variation on this loop.}

\newthought{Git has three areas:}

\begin{itemize}
    \item \emph{Working directory}\index{Git!working directory} The files on disk. What you edit.
    \item \emph{Staging area}\index{Git!staging area} What you have marked for the next commit. A deliberate selection.
    \item \emph{Repository}\index{Git!repository} The committed history. Permanent, immutable snapshots.
\end{itemize}

\newthought{The core commands:}

\marginnote{Write commit messages in the imperative mood: ``Add classifier module'' not ``Added classifier module.''
The message completes the sentence ``This commit will \ldots''. Keep the first line under 72 characters.
For a gallery of how not to do it, visit \url{https://whatthecommit.com}.}
\begin{verbatim}
git status              # what changed? what is staged?
git add file.py         # stage a file for the next commit
git add .               # stage everything (use with care)
git commit -m "message" # commit the staged changes
git diff                # unstaged changes vs last commit
git diff --staged       # staged changes vs last commit
git log                 # show commit history
git log --oneline       # compact view: one line per commit
\end{verbatim}

\marginnote{\newthought{When to commit?}
Commit when you have completed one logical change: a bug fix, a new function, a config update.
If you struggle to write a short commit message, you probably changed too many things at once.
If you go hours without committing, you are accumulating risk.
Small, frequent commits make history useful; large, rare commits make it a wall of text.}

\marginnote{\texttt{(\textbackslash\quad/)}\\
\texttt{(O.o)}\\
\texttt{(>\quad<)} Bunny approves these changes.}

% -----------------------------------------------------------------------------
\section{Branching and Merging}\index{Git!branching}\index{Git!merging}

\newthought{Branches let you work on ideas in isolation.} The \texttt{main} branch is the stable baseline.
Feature branches diverge, evolve, and merge back when ready.

\begin{verbatim}
git branch                   # list branches
git branch feature-x         # create a new branch
git checkout feature-x       # switch to it
git checkout -b feature-y    # create and switch in one step
git branch -d feature-x      # delete after merge
\end{verbatim}

\marginnote{Delete branches after merging. Stale branches accumulate fast
and make \texttt{git branch} unreadable. If the branch is merged, it is safe to delete;
the history lives on in \texttt{main}.}

\newthought{When the feature is ready,} merge it back and clean up:

\begin{verbatim}
git checkout main
git merge feature-x
git branch -d feature-x
\end{verbatim}

% -----------------------------------------------------------------------------
\subsection{Merge Conflicts}\index{Git!merge conflict}

\newthought{A merge conflict} occurs when two branches changed the same line in the same file.
Git cannot decide which version to keep, so you must resolve it manually.

\marginnote{Edit the file to keep the correct version and remove the markers.
Then stage the resolved file with \texttt{git add file.py}
and commit with \texttt{git commit -m "Resolve merge conflict in file.py"}.}

\newthought{Git marks the conflict} in the file. The block between \texttt{<<<<<<< HEAD} and \texttt{=======}
is what the current branch has; \texttt{HEAD} always means ``the commit you are on right now.''
The block between \texttt{=======} and \texttt{>>>>>>> feature-x} is what the incoming branch has.

\begin{verbatim}
<<<<<<< HEAD
result = model.predict(x)
=======
result = classifier.classify(x)
>>>>>>> feature-x
\end{verbatim}

% -----------------------------------------------------------------------------
\section{Tags}\index{Git!tags}

\newthought{A tag marks a point in history that matters:} a release, a milestone, a known-good state.

\begin{verbatim}
git tag v0.1                          # lightweight tag
git tag -a v0.1 -m "initial setup"    # annotated tag (preferred)
git push --tags                       # push tags to remote
\end{verbatim}

\marginnote{Lightweight tags are just a name pointing to a commit. Annotated tags store the tagger,
date, and message; use them for releases. We tag \texttt{v0.1} today; the tag becomes meaningful
when you look back at it in Block~8.\\[0.5em]
\texttt{tig}\index{tig} is a terminal-based Git browser that makes it easy to navigate commits, branches, and tags
on a headless Ubuntu machine. Install with \texttt{sudo apt install tig}.\\
\url{https://jonas.github.io/tig/}}

\newthought{Semantic versioning}\index{semantic versioning} gives tag names a shared meaning.
A version number follows the pattern \texttt{MAJOR.MINOR.PATCH}.
Increment \texttt{PATCH} for bug fixes that do not change behaviour,
\texttt{MINOR} for new features that remain backwards-compatible,
and \texttt{MAJOR} when you introduce breaking changes.
A jump from \texttt{v1.2.3} to \texttt{v1.3.0} tells users
that something was added but nothing they depend on was removed.
A jump to \texttt{v2.0.0} warns them to check for incompatibilities.

\marginnote{Semantic Versioning: \url{https://semver.org}}

% -----------------------------------------------------------------------------
\section{Remotes and GitHub}\index{Git!remotes}\index{GitHub}




\begin{verbatim}
    git clone <url>              # copy a remote repo locally
    # fork: copy repo to your GitHub account (GitHub UI, not a git command)
    git remote add origin git@github.com:user/pixelwise.git
\end{verbatim}

\newthought{A remote} is a copy of your repository on another machine, usually GitHub.

\newthought{Clone and fork}\index{Git!clone}\index{Git!fork} are two ways to 
start from an existing repository.

A \texttt{fork}\index{fork} is a GitHub concept, not a Git command.
It creates a full copy of someone else's repository under your own GitHub account.
You clone your fork locally, work on it freely, and propose changes back to the original via a pull request.
Forking is the standard workflow for contributing to open-source projects where you do not have write access to the original repository.

\marginnote{GitHub is not Git. Git is the version control system; GitHub is a hosting platform
that adds pull requests\index{pull request}, issues\index{issues}, CI/CD, and collaboration features on top of Git.
Alternatives include GitLab and Bitbucket.\\
GitHub: \url{https://github.com/}\\
GitHub Docs: \url{https://docs.github.com/}}

\texttt{git clone} copies a remote repository to your local machine, keeping the original as the \texttt{origin} remote.
You can pull updates and, if you have permission, push changes back.








% -----------------------------------------------------------------------------
\subsection{Pull Requests}

\newthought{Push}\index{Git!push} sends your local commits to the remote.
Until you push, your work exists only on your machine.

\newthought{Fetch}\index{Git!fetch} does the opposite: it downloads new commits from the remote
but leaves your working directory untouched, so you can inspect what changed before merging.

\newthought{Pull}\index{Git!pull} combines fetch and merge in one step,
updating your local branch to include the remote's latest commits.

\begin{verbatim}
git push -u origin main     # push and set upstream
git pull                     # fetch + merge from remote
git fetch                    # fetch without merging
\end{verbatim}

\newthought{The \texttt{-u} flag}\index{Git!-u flag} links your local branch to the remote branch.
After running \texttt{git push -u origin main} once, Git remembers the connection
and future \texttt{git push} and \texttt{git pull} work without specifying
the remote or branch.

\newthought{A pull request\index{pull request}} is a proposal to merge a branch into another branch, typically a
feature branch into \texttt{main}. It is the standard mechanism for code review: your changes are visible,
reviewable, and discussable before they become part of the main codebase.

\marginnote{\newthought{Pull requests and CI.}\index{continuous integration}
Because a PR represents a well-defined set of changes that has not yet reached \texttt{main},
it is the ideal moment to run a test suite, a linter, or any other check automatically.
If the tests fail, the merge is blocked and \texttt{main} stays healthy.
We will wire up this automation in a later block; for now, recognise that the PR
is where human review and machine verification meet.}

% -----------------------------------------------------------------------------
\newpage
\subsection{The \texttt{.gitignore} File}\index{.gitignore}

\newthought{Not everything belongs in version control.} The \texttt{.gitignore} file tells Git which files
and directories to exclude from tracking:

\begin{verbatim}
# Python
__pycache__/
*.pyc
.venv/

# Secrets
.env

# Data and models (too large, reproducible)
data/
models/*.pkl

# Editor files
.vscode/
*.swp
\end{verbatim}

\newthought{This includes all credentials:}\index{secrets} API keys, database passwords, webhook secrets,
tokens for AI services. These live on the server in environment variables or protected
configuration files, never in the repository.
A leaked production database password or API key can be exploited within minutes by automated scanners.

\marginnote{Git history is permanent\index{Git!history}. A secret committed once is leaked forever, even if you
delete the file in a later commit. The old commit still contains it. \texttt{.gitignore} is your first line of defence.}

% ==============================================================================
% SESSION B: HANDS-ON
% ==============================================================================

\newpage
% -----------------------------------------------------------------------------
\section{Exercise: Setting Up Git and the Repository}\index{exercises}

\marginnote{\newthought{Exercises} are for practice and reinforcing concepts.
Work through them yourself first, break things, discuss, this is not a time trial.
If your VM ends up in a state you cannot recover, restore the \texttt{B1-complete} snapshot and start over.
At the end of the session, take a snapshot named \texttt{B2-complete}, your safety net for the next block.}

\newthought{Create a GitHub account} if you do not have one. Then configure your Git identity on the \texttt{dev} VM:

\begin{verbatim}
git config --global user.name "Your Name"
git config --global user.email "your@email.com"
\end{verbatim}

Generate an SSH key and add it to GitHub so you can push over SSH.

\newthought{Initialise the PixelWise repo} on the \texttt{dev} VM:

\begin{verbatim}
mkdir pixelwise && cd pixelwise
git init
\end{verbatim}

Write a \texttt{README.md} with a one-line project description.


% -----------------------------------------------------------------------------
\section{Exercise: Staging, Committing, and .gitignore}\index{exercises}

\newthought{Write a \texttt{.gitignore}.} Create a \texttt{.gitignore} file that excludes
Python caches, virtual environments, secret files, dataset directories, model checkpoints, and editor files.
Verify it works: create a file called \texttt{scratch.pyc}, run \texttt{git status}, and confirm Git does not list it.
Then remove the test file.

\newthought{Make your first commit.} Stage and commit \texttt{README.md}
and \texttt{.gitignore}:

\begin{verbatim}
git add README.md .gitignore
git commit -m "Add project README and .gitignore"
\end{verbatim}

Run \texttt{git log} to see your first commit. Run \texttt{git status} to confirm the working directory is clean.

\marginnote{\newthought{The habit to build:} after every commit, run \texttt{git status} and \texttt{git log --oneline}
to confirm what just happened. It takes two seconds and prevents surprises.}


% -----------------------------------------------------------------------------
\newpage
\section{Exercise: Branching, Merging, and Conflicts}\index{exercises}

\newthought{Create a feature branch} and make a change to \texttt{README.md}:

\begin{verbatim}
git checkout -b feature-readme
# edit README.md: add a project description line
git add README.md
git commit -m "Add project description to README"
\end{verbatim}

\newthought{Switch back to \texttt{main}} and edit the same line differently:

\begin{verbatim}
git checkout main
# edit README.md: add a different description
#   on the same line
git add README.md
git commit -m "Add alternative description to README"
\end{verbatim}

\newthought{Merge and resolve the conflict.} Now merge the feature branch and observe the conflict:

\begin{verbatim}
git merge feature-readme
\end{verbatim}

Git will report a conflict. Open \texttt{README.md}, find the conflict markers
(\texttt{<<<<<<< HEAD}, \texttt{=======}, \texttt{>>>>>>>}), choose the correct version,
remove the markers, then complete the merge:

\begin{verbatim}
git add README.md
git commit -m "Resolve merge conflict in README"
\end{verbatim}

Verify with \texttt{git log --oneline --graph} that both branches are now integrated.


% -----------------------------------------------------------------------------
\section{Exercise: Tags and History Browsing}\index{exercises}

\newthought{Tag the release.} Mark this point in history:

\begin{verbatim}
git tag -a v0.1 -m "initial server setup"
git push --tags
\end{verbatim}

Verify the tag exists with \texttt{git tag -l} and inspect it with \texttt{git show v0.1}.
Confirm the tag appears on GitHub under the repository's ``Releases'' tab.

\newthought{Install \texttt{tig},} a terminal-based Git browser:

\begin{verbatim}
sudo apt install -y tig
tig
\end{verbatim}

Navigate the commit history and find your tag. Press \texttt{q} to quit.

\marginnote{\newthought{VS Code} can also browse Git history visually.
Open the PixelWise folder via \texttt{code .} or Remote SSH,
then use the Source Control panel and the GitLens extension
to explore commits, branches, and tags with a graphical interface.}


% -----------------------------------------------------------------------------
\newpage
\section{Exercise: Remotes, Push, and Clone}\index{exercises}

\newthought{Connect to GitHub.} Create a repository on GitHub, add it as a remote, and push:

\begin{verbatim}
git remote add origin git@github.com:user/pixelwise.git
git push -u origin main
\end{verbatim}

Verify on GitHub that your commits, tags, and \texttt{.gitignore} are all visible.

\newthought{Clone to the server.} SSH into the server and clone PixelWise to its permanent home:

\begin{verbatim}
sudo mkdir -p /opt/pixelwise
sudo chown student:student /opt/pixelwise
git clone git@github.com:user/pixelwise.git /opt/pixelwise
\end{verbatim}

PixelWise now lives on both machines. Run \texttt{git log --oneline} on the server to confirm the full history arrived.

\newthought{Take a snapshot} of both VMs. Name them \texttt{B2-complete}.

% -----------------------------------------------------------------------------
\section{Self-Reflection and Recap}

\newthought{Reflection questions:}
\begin{itemize}
    \item Why does changing a single character in a file produce a completely different commit hash?
    \item What is the difference between \texttt{git add} and \texttt{git commit}, and why does Git separate the two steps?
    \item When would you use a branch instead of committing directly to \texttt{main}, and what happens when a merge produces a conflict?
    \item Why should you never commit \texttt{.env} files or API keys, and what is the role of pull requests in catching mistakes before they reach \texttt{main}?
    \item Why does PixelWise use SSH keys rather than passwords for Git operations across both VMs?
\end{itemize}

\newthought{Milestone:} Your project lives on GitHub with a README, a \texttt{.gitignore}, and a tagged release.
The server has the repo cloned at \texttt{/opt/pixelwise}. Both machines share the same history.
In the next block we turn the manual package installs into a reproducible setup script and tackle dependency management.
